% GENERATED FILE. Edit canonical lessons, then run npm run generate:books.
% canonical-source-hash: fnv1a64:f474d55d265a5f2f
% canonical-lessons: IT-C470-piovigginare, IT-C470-schiarirsi, IT-C470-pozzanghera, IT-C470-fradicio, IT-C470-impermeabile, IT-R470-ripasso-469-470, IT-R470-ripasso-final-470

\chapter{Dopo la pioggia}
\label{ch:it-dopo-la-pioggia}
\hlchaptermodality{470}

\begin{chapteropening}
\textbf{By the end of this chapter:} \emph{I can talk about rainy weather and what follows: drizzle, the sky clearing, puddles, getting soaked and the raincoat.}

Complete the last lesson of chapter 470: I can talk about rainy weather and what follows: drizzle, the sky clearing, puddles, getting soaked and the raincoat.
\end{chapteropening}

\section[piovigginare]{piovigginare — to drizzle}

\label{lesson:IT-C470-piovigginare}



\begin{quote}
Before the new one: say the Italian for to blow, then the Italian for the cloudburst.
\end{quote}

\subsection*{You'll want to know: piovigginare}
\textbf{piovigginare} — \textquotedblleft{}to drizzle\textquotedblright{}.

\begin{quote}
Pioviggina, ma non serve l'ombrello.
\end{quote}

\begin{quote}
\emph{It is drizzling, but there is no need for an umbrella.}
\end{quote}

\emph{La pioggerella} — drizzle.

\begin{cousinweb}
\emph{piovigginare} is built on \emph{piovere} 'to rain' (EtymDB), which is inherited from Vulgar Latin \emph{plovere}, from Latin \emph{pluere}.
\end{cousinweb}

\begin{grammarlens}[title={What you've built}]
One more word for telling, explaining and describing at B1.
\end{grammarlens}

\subsection*{Your turn}
\begin{itemize}
\raggedright
  \item \emph{Say it:} \emph{piovigginare}
  \item \emph{Say it:} \emph{piovigginare}, once more
  \item \emph{Say it:} a sentence of your own with \emph{piovigginare}
  \item \emph{Recall:} say the Italian for to blow, then the Italian for the cloudburst, then say \emph{piovigginare} again
\end{itemize}

\begin{tcolorbox}[breakable,colback=teal!4,colframe=teal!35!black,title={Before you move on}]
What does piovigginare mean? (\textquotedblleft{}To drizzle\textquotedblright{}.) Say it once more.
\end{tcolorbox}
\section[schiarirsi]{schiarirsi — to clear up, to brighten}

\label{lesson:IT-C470-schiarirsi}



\begin{quote}
Before the new one: say the Italian for the cloudburst, then the Italian for to drizzle.
\end{quote}

\subsection*{You'll want to know: schiarirsi}
\textbf{schiarirsi} — \textquotedblleft{}to clear up, to brighten\textquotedblright{}.

\begin{quote}
Nel pomeriggio il cielo si schiarisce.
\end{quote}

\begin{quote}
\emph{In the afternoon the sky clears up.}
\end{quote}

\emph{Schiarirsi la voce} — to clear your throat.

\begin{cousinweb}
\emph{schiarirsi} is built from \emph{s-} and \emph{chiaro} 'clear', which is inherited from Latin \emph{clārus} (EtymDB).
\end{cousinweb}

\begin{grammarlens}[title={What you've built}]
One more word for telling, explaining and describing at B1.
\end{grammarlens}

\subsection*{Your turn}
\begin{itemize}
\raggedright
  \item \emph{Say it:} \emph{schiarirsi}
  \item \emph{Say it:} \emph{schiarirsi}, once more
  \item \emph{Say it:} a sentence of your own with \emph{schiarirsi}
  \item \emph{Recall:} say the Italian for the cloudburst, then the Italian for to drizzle, then say \emph{schiarirsi} again
\end{itemize}

\begin{tcolorbox}[breakable,colback=teal!4,colframe=teal!35!black,title={Before you move on}]
What does schiarirsi mean? (\textquotedblleft{}To clear up, to brighten\textquotedblright{}.) Say it once more.
\end{tcolorbox}
\section[la pozzanghera]{la pozzanghera — the puddle}

\label{lesson:IT-C470-pozzanghera}



\begin{quote}
Before the new one: say the Italian for to drizzle, then the Italian for to clear up.
\end{quote}

\subsection*{You'll want to know: la pozzanghera}
\textbf{la pozzanghera} — \textquotedblleft{}the puddle\textquotedblright{}.

\begin{quote}
Il bambino salta nelle pozzanghere.
\end{quote}

\begin{quote}
\emph{The child jumps in the puddles.}
\end{quote}

Built on \emph{la pozza}, a pool.

\begin{cousinweb}
\emph{pozzanghera} is built on \emph{pozza} 'pool', from \emph{pozzo} 'well', which is inherited from Latin \emph{puteus} (EtymDB).
\end{cousinweb}

\begin{grammarlens}[title={What you've built}]
One more word for telling, explaining and describing at B1.
\end{grammarlens}

\subsection*{Your turn}
\begin{itemize}
\raggedright
  \item \emph{Say it:} \emph{la pozzanghera}
  \item \emph{Say it:} \emph{la pozzanghera}, once more
  \item \emph{Say it:} a sentence of your own with \emph{la pozzanghera}
  \item \emph{Recall:} say the Italian for to drizzle, then the Italian for to clear up, then say \emph{la pozzanghera} again
\end{itemize}

\begin{tcolorbox}[breakable,colback=teal!4,colframe=teal!35!black,title={Before you move on}]
What does la pozzanghera mean? (\textquotedblleft{}The puddle\textquotedblright{}.) Say it once more.
\end{tcolorbox}
\section[fradicio]{fradicio — soaked, drenched}

\label{lesson:IT-C470-fradicio}



\begin{quote}
Before the new one: say the Italian for to clear up, then the Italian for the puddle.
\end{quote}

\subsection*{You'll want to know: fradicio}
\textbf{fradicio} — \textquotedblleft{}soaked, drenched\textquotedblright{}.

\begin{quote}
Sono arrivato a casa fradicio.
\end{quote}

\begin{quote}
\emph{I got home soaked.}
\end{quote}

\emph{Bagnato fradicio} — soaking wet; also rotten: \emph{una mela fradicia}.

\begin{cousinweb}
\emph{fradicio} is usually traced to Latin \emph{fracidus} 'soft, over-ripe' (Lewis \& Short), but the origin is not settled.
\end{cousinweb}

\begin{grammarlens}[title={What you've built}]
One more word for telling, explaining and describing at B1.
\end{grammarlens}

\subsection*{Your turn}
\begin{itemize}
\raggedright
  \item \emph{Say it:} \emph{fradicio}
  \item \emph{Say it:} \emph{fradicio}, once more
  \item \emph{Say it:} a sentence of your own with \emph{fradicio}
  \item \emph{Recall:} say the Italian for to clear up, then the Italian for the puddle, then say \emph{fradicio} again
\end{itemize}

\begin{tcolorbox}[breakable,colback=teal!4,colframe=teal!35!black,title={Before you move on}]
What does fradicio mean? (\textquotedblleft{}Soaked, drenched\textquotedblright{}.) Say it once more.
\end{tcolorbox}
\section[l'impermeabile]{l'impermeabile — the raincoat; waterproof}

\label{lesson:IT-C470-impermeabile}



\begin{quote}
Before the new one: say the Italian for the puddle, then the Italian for soaked.
\end{quote}

\subsection*{You'll want to know: l'impermeabile}
\textbf{l'impermeabile} — \textquotedblleft{}the raincoat; waterproof\textquotedblright{}.

\begin{quote}
Metti l'impermeabile, piove ancora.
\end{quote}

\begin{quote}
\emph{Put on your raincoat, it is still raining.}
\end{quote}

As an adjective: \emph{una giacca impermeabile}.

\begin{cousinweb}
\emph{impermeabile} is a learned borrowing from Late Latin \emph{impermeābilis} 'that cannot be passed through', built on \emph{meāre} 'to pass'.
\end{cousinweb}

\begin{grammarlens}[title={What you've built}]
One more word for telling, explaining and describing at B1.
\end{grammarlens}

\subsection*{Your turn}
\begin{itemize}
\raggedright
  \item \emph{Say it:} \emph{l'impermeabile}
  \item \emph{Say it:} \emph{l'impermeabile}, once more
  \item \emph{Say it:} a sentence of your own with \emph{l'impermeabile}
  \item \emph{Recall:} say the Italian for the puddle, then the Italian for soaked, then say \emph{l'impermeabile} again
\end{itemize}

\begin{tcolorbox}[breakable,colback=teal!4,colframe=teal!35!black,title={Before you move on}]
What does l'impermeabile mean? (\textquotedblleft{}The raincoat; waterproof\textquotedblright{}.) Say it once more.
\end{tcolorbox}
\section[(dialogue)]{Città, negozi, natura e tempo}

\label{lesson:IT-R470-ripasso-469-470}



\begin{quote}
Say the Italian for the blizzard and the hurricane.
\end{quote}

\subsection*{Your turn}
Say these aloud:

\begin{itemize}
\raggedright
  \item the blizzard, the hurricane, the gust, to blow, the cloudburst
  \item to drizzle, to clear up, the puddle, soaked, the raincoat
\end{itemize}

\begin{tcolorbox}[breakable,colback=teal!4,colframe=teal!35!black,title={Before you move on}]
The blizzard? (\textbf{la bufera}.) The hurricane? (\textbf{l'uragano}.) The gust? (\textbf{la raffica}.) To blow? (\textbf{soffiare}.) The cloudburst? (\textbf{il nubifragio}.) To drizzle? (\textbf{piovigginare}.) To clear up? (\textbf{schiarirsi}.) The puddle? (\textbf{la pozzanghera}.) Soaked? (\textbf{fradicio}.) The raincoat? (\textbf{l'impermeabile}.)
\end{tcolorbox}
\section[(dialogue)]{Dopo la pioggia — once more}

\label{lesson:IT-R470-ripasso-final-470}



\begin{quote}
Say the Italian for to drizzle and to clear up.
\end{quote}

\subsection*{Your turn}
\begin{itemize}
\raggedright
  \item \emph{Say it:} to drizzle, to clear up, the puddle, soaked, the raincoat
\end{itemize}

\begin{tcolorbox}[breakable,colback=teal!4,colframe=teal!35!black,title={Before you move on}]
To drizzle? (\textbf{piovigginare}.) To clear up? (\textbf{schiarirsi}.) The puddle? (\textbf{la pozzanghera}.) Soaked? (\textbf{fradicio}.) The raincoat? (\textbf{l'impermeabile}.)
\end{tcolorbox}
